\documentclass{article}
\usepackage[utf8]{inputenc}
\usepackage{polski}
\usepackage{classdiagram}


\usepackage[a4paper, margin=1.5cm]{geometry}
\usepackage{tikz}
\usetikzlibrary{calc}
\usepackage{pdfpages}

\makeatletter
\def\firstauthor#1{\gdef\@firstauthor{#1}}
\def\@firstauthor{\@latex@warning@no@line{No \noexpand\firstauthor given}}
\def\secondauthor#1{\gdef\@secondauthor{#1}}
\def\@secondauthor{}
\def\thirdauthor#1{\gdef\@thirdauthor{#1}}
\def\@thirdauthor{}
\def\projectname#1{\gdef\@projectname{#1}}
\def\@projectname{\@latex@warning@no@line{No \noexpand\projectname given}}
\newcommand\makeheader{%

\begin{tikzpicture}
    \node(IN){Imię i nazwisko};
    \node(IN1) at ($(IN) + (0, -0.5)$) {\@firstauthor};
    \node(IN2) at ($(IN1) + (0, -0.5)$) {\@secondauthor};
    \node(IN3) at ($(IN2) + (0, -0.5)$) {\@thirdauthor};
\end{tikzpicture}
\vspace*{20pt}
\begin{center}
    \begin{tikzpicture}
        \node(PN){\bf\LARGE\@projectname};
        \node(DOK) at ($(PN) + (0, 1)$) {\Large DOKUMENTACJA TECHNICZNA PROJEKTU};
        \node(PDM) at ($(PN) + (0, -1)$) {\Large Przedmiot: PROGRAMOWANIE ZAAWANSOWANE};
    \end{tikzpicture}
\end{center}
}
\makeatother


\projectname{[Webowa Lista Zadań Zespołu]}
\firstauthor{[Aleksandra Chmielewska]}
\secondauthor{[Wojciech Domżalski]}
\thirdauthor{[Eryk Dolajzer]}

\begin{document}

\makeheader

\section{Tematyka projektu}
{\small\it Projekt polega na stworzeniu narzędzia zdolnego do przechowywania aktualnych zadań dla zespołu lub osoby. Ma być maksymalnie prosty i funkcjonalny, umożliwiający szybkie zorientowanie się w aktywnych zadaniach.}

\section{Diagram klas}
\includegraphics[scale=0.50] {diagram.pdf}

\section{Wymagania systemowe}
\subsection{Wymagania funkcjonalne}
\begin{enumerate}
    \item Dodawanie Zadań
    \item Usuwanie Zadań
    \item Przeglądanie Zadań
    \item Przechowywanie Zadań
    \item Opisywanie Zadań
    \item Numerowanie Zadań
\end{enumerate}
\subsection{Wymagania pozafunkcjonalne}
\begin{enumerate}
    \item Dostęp za pomoca przeglądarki
    \item Prosty interfejs
    \item Kolejność FIFO
\end{enumerate}
\section{Opis techniczny projektu}
{\small\it 
\begin{enumerate}
    \item Realizacja za pomocą wzorca projektowego MVC
    \item Realizacja za pomocą bibliotek natywnych ASP .NET 5.0
    \item Realizacja na frameworku .NET 5.0
    \item Realizacja z wykorzystaniem Bazy Danych MS.SQL
    \item Interakcja z użytkownikiem za pomocą warstwy VIEW, z wykorzystaniem CSHTML
\end{enumerate}
\section{Potencjalne możliwe problemy i zagrożenia}
{\small\it Projekt nie posiada błędów które są mozliwe do okreslenie na etapie planowania.}
\section{Scenariusze testów}
{\small\it Procedura Testowa}
\begin{enumerate}
    \item Uruchomienie
    \item Wprowadzenie nowego zadania
    \item Próba przekroczenia dostępnej pamięci dla kontenerów
    \item Akceptacja Zadania
    \item Edytowanie Zadania
     \item   Próba przekroczenia dostępnej pamięci dla kontenerów
     \item Akceptacja zadania
     \item Sprawdzenie szczegółów zadania
     \item Próba usunięcia zadania
     \item Dodanie zadania
     \item wyłączenie programu
     \item Uruchomienie programu
     \item Sprawdzenie istnienia zadania
     \item Sprawdzenie poprawności zawartości zadania
     \item koniec testów
\end{enumerate}
\pagebreak
\section{Lista zmian w dokumencie}

\begin{table}[h!]
\centering
 \begin{tabular}{|c c c |} 
 \hline
 Rewizja & Imię i nazwisko & Opis \\ [0.5ex] 
 \hline\hline
 1 & Jacek Kuźmicz & Stworzenie szablonu dokumentu \\ 
 2 & Wojciech Frohmberg & Poprawki i wersja LaTeX  \\
 3 & Wojciech Domżalski & Opisanie ogólnych założeń \\
 4 & Eryk Dolajzer  &   Wymagania systemowe i funkcjonalne \\
 5 & Aleksandra Chmielewska & Scenariusze testów i Wymagania pozafunkcjonalne \\
 6 & Wojciech Domżalski & Diagram Klas i opis techniczny projektu \\
 \hline
 \end{tabular}
\end{table}

\end{document}

